\tikzset{
		axpic/.style={baseline=(current bounding box.center), x=0.6cm, y=0.52cm, line cap=round},
		boxnode/.style={rounded rectangle, fill=white, draw=black, inner sep=2pt, anchor=center}
	}

	\medskip
	\noindent\fbox{\begin{minipage}{0.97\linewidth}
			\begin{center}\textbf{$\mathbf{GLA}$ axioms}\quad{\footnotesize\cite{paixao2022highlevel,bonchi2017interacting}}\end{center}

			\textit{White commutative monoid (addition $\oplus$, unit $\circ$).}
			\begin{center}\begin{tabular}{r@{\quad}c@{\;$=$\;}c}
					(assoc) &
                    \begin{tikzpicture}
                      \begin{scope}[shift={(0.000,1.000)}]
                      \pic (g1) at (0.500,0.500) {multiplication};
                      \end{scope}
                      \coordinate (id_in_2) at (0,0.500);
                      \coordinate (id_out_3) at (1.000,0.500);
                      \draw (id_in_2) -- (id_out_3);
                      \begin{scope}[shift={(1.250,0.500)}]
                      \pic (g4) at (0.500,0.500) {multiplication};
                      \end{scope}
                      \draw (g1-out-0) to[out=0,in=180] (g4-in-0);
                      \draw (id_out_3) to[out=0,in=180] (g4-in-1);
                    \end{tikzpicture}
					        &
                    \begin{tikzpicture}
                      \begin{scope}[shift={(0.000,1.000)}]
                      \coordinate (id_in_1) at (0,0.500);
                      \coordinate (id_out_2) at (1.000,0.500);
                      \draw (id_in_1) -- (id_out_2);
                      \end{scope}
                      \pic (g3) at (0.500,0.500) {multiplication};
                      \begin{scope}[shift={(1.250,0.500)}]
                      \pic (g4) at (0.500,0.500) {multiplication};
                      \end{scope}
                      \draw (id_out_2) to[out=0,in=180] (g4-in-0);
                      \draw (g3-out-0) to[out=0,in=180] (g4-in-1);
                    \end{tikzpicture}
					\\[3mm]
					(comm)  &
                    \begin{tikzpicture}
                      \coordinate (sw_in_top_1) at (0,1.000);
                      \coordinate (sw_in_bottom_2) at (0,0.000);
                      \coordinate (sw_out_top_3) at (1.000,1.000);
                      \coordinate (sw_out_bottom_4) at (1.000,0.000);
                      \draw (sw_in_bottom_2) .. controls (0.500,0.000) and (0.500,1.000) .. (sw_out_top_3);
                      \draw (sw_in_top_1) .. controls (0.500,1.000) and (0.500,0.000) .. (sw_out_bottom_4);
                      \begin{scope}[shift={(1.250,0.000)}]
                      \pic (g5) at (0.500,0.500) {multiplication};
                      \end{scope}
                      \draw (sw_out_top_3) to[out=0,in=180] (g5-in-0);
                      \draw (sw_out_bottom_4) to[out=0,in=180] (g5-in-1);
                    \end{tikzpicture}
					        &
                    \begin{tikzpicture}
                      \pic (g1) at (0.500,0.500) {multiplication};
                    \end{tikzpicture}
					\\[3mm]
					(unit)  &
                    \begin{tikzpicture}
                      \begin{scope}[shift={(0.000,1.000)}]
                      \pic (g1) at (0.500,0.500) {zero};
                      \end{scope}
                      \coordinate (id_in_2) at (0,0.500);
                      \coordinate (id_out_3) at (1.000,0.500);
                      \draw (id_in_2) -- (id_out_3);
                      \begin{scope}[shift={(1.250,0.500)}]
                      \pic (g4) at (0.500,0.500) {multiplication};
                      \end{scope}
                      \draw (g1-out-0) to[out=0,in=180] (g4-in-0);
                      \draw (id_out_3) to[out=0,in=180] (g4-in-1);
                    \end{tikzpicture}
					        &
                    \begin{tikzpicture}
                      \coordinate (id_in_1) at (0,0.500);
                      \coordinate (id_out_2) at (1.000,0.500);
                      \draw (id_in_1) -- (id_out_2);
                    \end{tikzpicture}
				\end{tabular}\end{center}

			\textit{Black cocommutative comonoid (copy, counit $\bullet$).}
			\begin{center}\begin{tabular}{r@{\quad}c@{\;$=$\;}c}
					(coassoc) &
                    \begin{tikzpicture}
                      \pic (g1) at (0.500,0.500) {copy};
                      \begin{scope}[shift={(1.250,0.000)}]
                      \pic (g2) at (0.500,0.500) {copy};
                      \end{scope}
                      \coordinate (id_in_3) at (1.250,1.000);
                      \coordinate (id_out_4) at (2.250,1.000);
                      \draw (id_in_3) -- (id_out_4);
                      \draw (g1-out-0) to[out=0,in=180] (g2-in-0);
                      \draw (g1-out-1) to[out=0,in=180] (id_in_3);
                    \end{tikzpicture}
					          &
                    \begin{tikzpicture}
                      \begin{scope}[shift={(0.000,1.000)}]
                      \pic (g1) at (0.500,0.500) {copy};
                      \end{scope}
                      \begin{scope}[shift={(1.250,0.000)}]
                      \pic (g2) at (0.500,0.500) {copy};
                      \end{scope}
                      \coordinate (id_in_3) at (1.250,0.000);
                      \coordinate (id_out_4) at (2.250,0.000);
                      \draw (id_in_3) -- (id_out_4);
                      \draw (g1-out-1) to[out=0,in=180] (g2-in-0);
                      \draw (g1-out-0) to[out=0,in=180] (id_in_3);
                    \end{tikzpicture}
					\\[3mm]
					(cocomm)  &
                    \begin{tikzpicture}
                      \pic (g1) at (0.500,0.500) {copy};
                      \coordinate (sw_in_top_1) at (1.000,1.000);
                      \coordinate (sw_in_bottom_2) at (1.000,0.000);
                      \coordinate (sw_out_top_3) at (2.000,1.000);
                      \coordinate (sw_out_bottom_4) at (2.000,0.000);
                      \draw (g1-out-1) to[out=0,in=180] (sw_in_top_1);
                      \draw (g1-out-0) to[out=0,in=180] (sw_in_bottom_2);
                      \draw (sw_in_bottom_2) .. controls (1.500,0.000) and (1.500,1.000) .. (sw_out_top_3);
                      \draw (sw_in_top_1) .. controls (1.500,1.000) and (1.500,0.000) .. (sw_out_bottom_4);
                    \end{tikzpicture}
					          &
                    \begin{tikzpicture}
                      \pic (g1) at (0.500,0.500) {copy};
                    \end{tikzpicture}
					\\[3mm]
					(counit)  &
                    \begin{tikzpicture}
                      \pic (g1) at (0.500,0.500) {copy};
                      \begin{scope}[shift={(1.250,0.500)}]
                      \pic (g2) at (0.500,0.500) {discard};
                      \end{scope}
                      \coordinate (id_in_3) at (1.250,0.000);
                      \coordinate (id_out_4) at (2.250,0.000);
                      \draw (id_in_3) -- (id_out_4);
                      \draw (g1-out-1) to[out=0,in=180] (g2-in-0);
                      \draw (g1-out-0) to[out=0,in=180] (id_in_3);
                    \end{tikzpicture}
					          &
                    \begin{tikzpicture}
                      \coordinate (id_in_1) at (0,0.500);
                      \coordinate (id_out_2) at (1.000,0.500);
                      \draw (id_in_1) -- (id_out_2);
                    \end{tikzpicture}
				\end{tabular}\end{center}

			\textit{Bialgebra / Hopf (white over black).}
			\begin{center}\begin{tabular}{r@{\quad}c@{\;$=$\;}c}
					(B1)       &
                    \begin{tikzpicture}
                      \pic (g1) at (0.500,0.500) {multiplication};
                      \begin{scope}[shift={(1.250,0.000)}]
                      \pic (g2) at (0.500,0.500) {copy};
                      \end{scope}
                      \draw (g1-out-0) to[out=0,in=180] (g2-in-0);
                    \end{tikzpicture}
					           &
                    \begin{tikzpicture}
                      \begin{scope}[shift={(0.000,1.500)}]
                      \pic (g1) at (0.500,0.500) {copy};
                      \end{scope}
                      \begin{scope}[shift={(0.000,0.000)}]
                      \pic (g2) at (0.500,0.500) {copy};
                      \end{scope}
                      \begin{scope}[shift={(1.750,1.500)}]
                      \pic (g3) at (0.500,0.500) {multiplication};
                      \end{scope}
                      \begin{scope}[shift={(1.750,0.000)}]
                      \pic (g4) at (0.500,0.500) {multiplication};
                      \end{scope}
                      \draw (g1-out-1) to[out=0,in=180] (g3-in-0);
                      \draw (g2-out-0) to[out=0,in=180] (g4-in-1);
                      \draw (g1-out-0) to[out=0,in=180] (g4-in-0);
                      \draw (g2-out-1) to[out=0,in=180] (g3-in-1);
                    \end{tikzpicture}
					\\[3mm]
					(B2)       &
                    \begin{tikzpicture}
                      \pic (g1) at (0.500,0.500) {multiplication};
                      \begin{scope}[shift={(1.250,0.500)}]
                      \pic (g2) at (0.500,0.500) {discard};
                      \end{scope}
                      \draw (g1-out-0) to[out=0,in=180] (g2-in-0);
                    \end{tikzpicture}
					           &
                    \begin{tikzpicture}
                      \begin{scope}[shift={(0.000,0.700)}]
                      \pic (g1) at (0.500,0.500) {discard};
                      \end{scope}
                      \begin{scope}[shift={(0.000,-0.700)}]
                      \pic (g2) at (0.500,0.500) {discard};
                      \end{scope}
                    \end{tikzpicture}
					\\[3mm]
					(B3)       &
                    \begin{tikzpicture}
                      \pic (g1) at (0.500,0.500) {zero};
                      \begin{scope}[shift={(1.250,0.000)}]
                      \pic (g2) at (0.500,0.500) {copy};
                      \end{scope}
                      \draw (g1-out-0) to[out=0,in=180] (g2-in-0);
                    \end{tikzpicture}
					           &
                    \begin{tikzpicture}
                      \begin{scope}[shift={(0.000,0.700)}]
                      \pic (g1) at (0.500,0.500) {zero};
                      \end{scope}
                      \begin{scope}[shift={(0.000,-0.700)}]
                      \pic (g2) at (0.500,0.500) {zero};
                      \end{scope}
                    \end{tikzpicture}
					\\[3mm]
					(B4, unit) &
                    \begin{tikzpicture}
                      \pic (g1) at (0.500,0.500) {zero};
                      \begin{scope}[shift={(1.250,0.000)}]
                      \pic (g2) at (0.500,0.500) {discard};
                      \end{scope}
                      \draw (g1-out-0) to[out=0,in=180] (g2-in-0);
                    \end{tikzpicture}
					           &
                    \begin{tikzpicture}
                      \path (0,0) rectangle (1.000,1.000);
                    \end{tikzpicture}
				\end{tabular}\end{center}

			\textit{Frobenius and special law (copy structure; the addition structure is dual).}
			\begin{center}\begin{tabular}{r@{\quad}c@{\;$=$\;}c}
					(Frob)    &
                    \begin{tikzpicture}
                      \begin{scope}[shift={(0.000,1.000)}]
                      \pic (g1) at (0.500,0.500) {copy};
                      \end{scope}
                      \begin{scope}[shift={(1.250,0.000)}]
                      \pic (g2) at (0.500,0.500) {cocopy};
                      \end{scope}
                      \coordinate (id_in_3) at (1.000,1.500);
                      \coordinate (id_out_4) at (2.750,1.500);
                      \draw (id_in_3) -- (id_out_4);
                      \coordinate (id_in_5) at (1.000,0.000);
                      \coordinate (id_out_6) at (2.750,0.000);
                      \draw (id_in_5) -- (id_out_6);
                      \draw (g1-out-1) to[out=0,in=180] (g2-in-0);
                      \draw (g1-out-0) to[out=0,in=180] (id_in_3);
                      \draw (id_out_6) to[out=0,in=0] (g2-in-1);
                    \end{tikzpicture}
					          &
                    \begin{tikzpicture}
                      \pic (g1) at (0.500,0.500) {cocopy};
                      \begin{scope}[shift={(1.250,0.000)}]
                      \pic (g2) at (0.500,0.500) {copy};
                      \end{scope}
                      \draw (g1-out-0) to[out=0,in=180] (g2-in-0);
                    \end{tikzpicture}
					\\[3mm]
					(special) &
                    \begin{tikzpicture}
                      \pic (g1) at (0.500,0.500) {copy};
                      \begin{scope}[shift={(1.250,0.000)}]
                      \pic (g2) at (0.500,0.500) {cocopy};
                      \end{scope}
                      \draw (g1-out-0) to[out=0,in=180] (g2-in-0);
                      \draw (g1-out-1) to[out=0,in=180] (g2-in-1);
                    \end{tikzpicture}
					          &
                    \begin{tikzpicture}
                      \coordinate (id_in_1) at (0,0.500);
                      \coordinate (id_out_2) at (1.000,0.500);
                      \draw (id_in_1) -- (id_out_2);
                    \end{tikzpicture}
				\end{tabular}\end{center}

			\textit{Scalars and antipode ($\text{-}1$).}
			\begin{center}\begin{tabular}{r@{\quad}c@{\;$=$\;}c}
					(s.\,add)  &
                    \begin{tikzpicture}
                      \pic (g1) at (0.500,0.500) {copy};
                      \begin{scope}[shift={(1.250,0.700)}]
                      \pic [val=k] (g2) at (0.500,0.500) {scalar};
                      \end{scope}
                      \begin{scope}[shift={(1.250,-0.700)}]
                      \pic [val=l] (g3) at (0.500,0.500) {scalar};
                      \end{scope}
                      \begin{scope}[shift={(2.750,0.000)}]
                      \pic (g4) at (0.500,0.500) {multiplication};
                      \end{scope}
                      \draw (g1-out-1) to[out=0,in=180] (g2-in-0);
                      \draw (g1-out-0) to[out=0,in=180] (g3-in-0);
                      \draw (g2-out-0) to[out=0,in=180] (g4-in-0);
                      \draw (g3-out-0) to[out=0,in=180] (g4-in-1);
                    \end{tikzpicture}
					           &
                    \begin{tikzpicture}
                      \pic [val={k+l}] (g1) at (0.500,0.500) {scalar};
                    \end{tikzpicture}
					\\[3mm]
					(s.\,mult) &
                    \begin{tikzpicture}
                      \pic [val=k] (g1) at (0.500,0.500) {scalar};
                      \begin{scope}[shift={(1.250,0.000)}]
                      \pic [val=l] (g2) at (0.500,0.500) {scalar};
                      \end{scope}
                      \draw (g1-out-0) to[out=0,in=180] (g2-in-0);
                    \end{tikzpicture}
					           &
                    \begin{tikzpicture}
                      \pic [val={kl}] (g1) at (0.500,0.500) {scalar};
                    \end{tikzpicture}
					\\[3mm]
					(antipode) &
                    \begin{tikzpicture}
                      \pic (g1) at (0.500,0.500) {copy};
                      \begin{scope}[shift={(1.250,-0.700)}]
                      \pic [val={\text{-}1}] (g2) at (0.500,0.500) {scalar};
                      \end{scope}
                      \begin{scope}[shift={(2.750,0.000)}]
                      \pic (g3) at (0.500,0.500) {multiplication};
                      \end{scope}
                      \coordinate (id_in_4) at (1.250,1.200);
                      \coordinate (id_out_5) at (2.750,1.200);
                      \draw (id_in_4) -- (id_out_5);
                      \draw (g1-out-1) to[out=0,in=180] (id_in_4);
                      \draw (g1-out-0) to[out=0,in=180] (g2-in-0);
                      \draw (id_out_5) to[out=0,in=180] (g3-in-0);
                      \draw (g2-out-0) to[out=0,in=180] (g3-in-1);
                    \end{tikzpicture}
					           &
                    \begin{tikzpicture}
                      \pic (g1) at (0.500,0.500) {discard};
                      \begin{scope}[shift={(1.250,0.000)}]
                      \pic (g2) at (0.500,0.500) {zero};
                      \end{scope}
                    \end{tikzpicture}
				\end{tabular}\end{center}
		\end{minipage}}

	\medskip
	\noindent\fbox{\begin{minipage}{0.97\linewidth}
			\begin{center}\textbf{$\mathbf{GAA}$ axioms (in addition to $\mathbf{GLA}$)}\quad{\footnotesize\cite{bonchi2019affine}}\end{center}
			The affine generator $\rtop$ (a point $0\to 1$, drawn as the empty box) is a comonoid homomorphism:
			\begin{center}\begin{tabular}{r@{\quad}c@{\;$=$\;}c}
					(aff.\,copy) &
                    \begin{tikzpicture}
                      \pic (g1) at (0.500,0.500) {affine};
                      \begin{scope}[shift={(1.250,0.000)}]
                      \pic (g2) at (0.500,0.500) {copy};
                      \end{scope}
                      \draw (g1-out-0) to[out=0,in=180] (g2-in-0);
                    \end{tikzpicture}
					             &
                    \begin{tikzpicture}
                      \begin{scope}[shift={(0.000,0.700)}]
                      \pic (g1) at (0.500,0.500) {affine};
                      \end{scope}
                      \begin{scope}[shift={(0.000,-0.700)}]
                      \pic (g2) at (0.500,0.500) {affine};
                      \end{scope}
                    \end{tikzpicture}
					\\[3mm]
					(aff.\,del)  &
                    \begin{tikzpicture}
                      \pic (g1) at (0.500,0.500) {affine};
                      \begin{scope}[shift={(1.250,0.000)}]
                      \pic (g2) at (0.500,0.500) {discard};
                      \end{scope}
                      \draw (g1-out-0) to[out=0,in=180] (g2-in-0);
                    \end{tikzpicture}
					             &
                    \begin{tikzpicture}
                      \path (0,0) rectangle (1.000,1.000);
                    \end{tikzpicture}
				\end{tabular}\end{center}
		\end{minipage}}

	\medskip
	\noindent\fbox{
        \begin{minipage}{0.97\linewidth}
			\begin{center}\textbf{$\mathbf{GPA}$ axioms (in addition to $\mathbf{GAA}$)}\quad{\footnotesize\cite{bonchi2021polyhedral}}\end{center}
			The order primitive $\ge$ is a preorder compatible with the linear structure
			(transitivity; monoid and comonoid morphism; $0\ge 0$):
			\begin{center}\begin{tabular}{r@{\quad}c@{\;$=$\;}c}
					(trans)           &
                    \begin{tikzpicture}
                      \pic (g1) at (0.500,0.500) {ge};
                      \begin{scope}[shift={(1.250,0.000)}]
                      \pic (g2) at (0.500,0.500) {ge};
                      \end{scope}
                      \draw (g1-out-0) to[out=0,in=180] (g2-in-0);
                    \end{tikzpicture}
					                  &
                    \begin{tikzpicture}
                      \pic (g1) at (0.500,0.500) {ge};
                    \end{tikzpicture}
					\\[3mm]
					(mono.\,$\oplus$) &
                    \begin{tikzpicture}
                      \pic (g1) at (0.500,0.500) {multiplication};
                      \begin{scope}[shift={(1.250,0.000)}]
                      \pic (g2) at (0.500,0.500) {ge};
                      \end{scope}
                      \draw (g1-out-0) to[out=0,in=180] (g2-in-0);
                    \end{tikzpicture}
					                  &
                    \begin{tikzpicture}
                      \begin{scope}[shift={(0.000,0.800)}]
                      \pic (g1) at (0.500,0.500) {ge};
                      \end{scope}
                      \begin{scope}[shift={(0.000,-0.800)}]
                      \pic (g2) at (0.500,0.500) {ge};
                      \end{scope}
                      \begin{scope}[shift={(1.500,0.000)}]
                      \pic (g3) at (0.500,0.500) {multiplication};
                      \end{scope}
                      \draw (g1-out-0) to[out=0,in=180] (g3-in-0);
                      \draw (g2-out-0) to[out=0,in=180] (g3-in-1);
                    \end{tikzpicture}
					\\[3mm]
					(mono.\,copy)     &
                    \begin{tikzpicture}
                      \pic (g1) at (0.500,0.500) {copy};
                      \begin{scope}[shift={(1.250,0.800)}]
                      \pic (g2) at (0.500,0.500) {ge};
                      \end{scope}
                      \begin{scope}[shift={(1.250,-0.800)}]
                      \pic (g3) at (0.500,0.500) {ge};
                      \end{scope}
                      \draw (g1-out-1) to[out=0,in=180] (g2-in-0);
                      \draw (g1-out-0) to[out=0,in=180] (g3-in-0);
                    \end{tikzpicture}
					                  &
                    \begin{tikzpicture}
                      \pic (g1) at (0.500,0.500) {ge};
                      \begin{scope}[shift={(1.250,0.000)}]
                      \pic (g2) at (0.500,0.500) {copy};
                      \end{scope}
                      \draw (g1-out-0) to[out=0,in=180] (g2-in-0);
                    \end{tikzpicture}
					\\[3mm]
					(zero)            &
                    \begin{tikzpicture}
                      \pic (g1) at (0.500,0.500) {zero};
                      \begin{scope}[shift={(1.250,0.000)}]
                      \pic (g2) at (0.500,0.500) {ge};
                      \end{scope}
                      \draw (g1-out-0) to[out=0,in=180] (g2-in-0);
                    \end{tikzpicture}
					                  &
                    \begin{tikzpicture}
                      \pic (g1) at (0.500,0.500) {zero};
                    \end{tikzpicture}
				\end{tabular}\end{center}
			Together with reflexivity, antisymmetry and positivity of scalars; see \cite{bonchi2021polyhedral,bonchi2021.9farkas}.
		\end{minipage}
    }
